\label{cdp:st:C.2.UN.13}
\StatementRecord{cdp:st:C.2.UN.13}{proved}{}{the hypotheses, notation and definitions stated here; \cref{cdp:eq:C.KM.4}; \cref{cdp:eq:C.KM.5}; \cite[Chapter~1, Example~1.3(iii)]{HuybrechtsK3Draft}; \cref{cdp:st:C.2.UN.4}; \cref{cdp:eq:C.KM.S2}; \cref{cdp:st:C.2.UN.3}}

Fix an abelian surface $A$. There is a smooth projective threefold
$W$ with a proper map $f_W:W\to\PP^1$ whose fibres over
$\C^*$ are isomorphic to $A$. It has precisely two special fibres,
over $0$ and $\infty$. Each has full divisor
\[
               2P+\sum_{i=1}^{16}E_i,\qquad
 P=\operatorname{Km}(A),\quad E_i\simeq\PP^2.       \tag{C.KM.1}\label{cdp:eq:C.KM.1}
\]
Here $P$ is the smooth Kummer K3 surface, the $E_i$ are disjoint,
and $P\cap E_i$ is a projective line on $E_i$ and a curve of
self-intersection $-2$ on $P$. The reduced divisor is simple
normal crossing and has Euler number $40$.

There is a composition of smooth vertical blowups $Y\to W$,
disjoint from both K3 components, for which both special reduced
fibres have Euler number zero. Each has $57$ irreducible
components, including its original unchanged Kummer surface.
The maps to the base and the nearby tori are retained. Their
rational Betti numbers and coherent dimensions are
\[
 \begin{array}{c|c|c|c}
 & (b_i)_{i=0}^6 & e & (h^i(\OO))_{i=0}^3\\ \hline
 W&(1,0,39,0,39,0,1)&80&(1,0,1,0)\\
 Y&(1,0,119,240,119,0,1)&0&(1,0,1,0).
 \end{array}                                             \tag{C.KM.2}\label{cdp:eq:C.KM.2}
\]
For either special value, let $S_W$ and $S_Y$ denote its
reduced supports in these two models. Their cohomology is
\[
 \begin{gathered}
 (H^i(S_W,\Q))_{i=0}^4=(\Q,0,\Q^{22},0,\Q^{17}),\\
 (H^i(S_Y,\Q))_{i=0}^4=(\Q,0,\Q^{62},\Q^{120},\Q^{57}),\\
 (h^i(S_W,\OO_{S_W}))_{i=0}^2
   =(h^i(S_Y,\OO_{S_Y}))_{i=0}^2=(1,0,1).
 \end{gathered}                                          \tag{C.KM.S1}\label{cdp:eq:C.KM.S1}
\]
Neither K3 component can be contracted by a bimeromorphic map
between smooth fillings. For comparison, after the quadratic
base change the normalization of the original quotient is
exactly $A\times\PP^1$, and the corresponding base changes
of $W$ and $Y$ are birational modifications of that product.
Thus the asserted persistence concerns bimeromorphic comparison
over the original base; its relation to base change is given
by this specified map.
\end{proposition}


\begin{proof}
\noindent\textbf{Step 1.} On $A\times\PP^1_s$ take the involution
\[
       \iota(a,s)=(-a,-s),\qquad
 X=(A\times\PP^1_s)/\langle\iota\rangle,
       \qquad t=s^2.                                  \tag{C.KM.3}\label{cdp:eq:C.KM.3}
\]
The quotient is projective and normal. Its fixed points are
exactly $A[2]\times\{0,\infty\}$, so there are $32$ of them.
Away from them the action is free. In local coordinates centred
at any fixed point, including the coordinate $s^{-1}$ at
infinity, the action is scalar multiplication by $-1$ on
$\C^3$. Thus the only singularities of $X$ are those $32$
quotient points.

\noindent\textbf{Step 2.} Blow up these points in $X$, and call the resulting map
$\pi:W\to X$. We verify its smoothness and its exact relation
to the quotient before using it. Write the coordinates of a
covering chart as $(a,b,s)$. The invariant local ring is
generated by their six quadratic monomials. In the blowup
chart corresponding to $a^2$, put
\[
             a=r,\qquad b=ru,\qquad s=rv,
             \qquad R=r^2.                            \tag{C.KM.4}\label{cdp:eq:C.KM.4}
\]
All invariant generators pull back as
\[
 (a^2,ab,as,b^2,bs,s^2)
           =(R,Ru,Rv,Ru^2,Ruv,Rv^2).                  \tag{C.KM.5}\label{cdp:eq:C.KM.5}
\]
Consequently $(R,u,v)$ are smooth coordinates on the blowup.
The analogous charts for $b^2$ and $s^2$ cover it: the
exceptional projectivized invariant cone is the quadratic
Veronese image of $\PP^2$, on which at least one of the
three squared coordinates is nonzero. Equivalently, blow up
the origin of $\C^3$ first. The lifted involution fixes its
exceptional $\PP^2$ and sends its normal coordinate $r$ to
$-r$; its quotient has precisely \eqref{cdp:eq:C.KM.4}--\eqref{cdp:eq:C.KM.5}.
The transitions for $R=r^2$ identify its normal line with
$\OO_{\PP^2}(-2)$. This proves that $W$ is smooth and
projective, with exceptional divisors $E_i\simeq\PP^2$.

\noindent\textbf{Step 3.} The base map in \eqref{cdp:eq:C.KM.4} is exactly
\[
                             t=R v^2.                 \tag{C.KM.6}\label{cdp:eq:C.KM.6}
\]
Here $E_i=(R=0)$ and the strict central component is
$P=(v=0)$. In the chart where $s$ is the exceptional
coordinate the map is $t=R$ and only $E_i$ occurs.
These charts prove the multiplicities and the transverse
intersection in \eqref{cdp:eq:C.KM.1}. They also show that
$N_{E_i/W}|_{P\cap E_i}=\OO_{\PP^1}(-2)$, which is
the normal line of that intersection in $P$.
The same calculation applies over infinity. Off those two
values, choose either square root $s$ of $t$; the map
$a\mapsto[a,s]$ identifies the fibre with $A$, and switching
the square root composes this identification with $a\mapsto-a$.
In particular the map has connected fibres and its monodromy
on $H^1(A,\Z)$ is $-I_4$.

\noindent\textbf{Step 4.} The restriction $P\to A/\{\pm1\}$ resolves exactly the
sixteen quotient surface nodes. It is the classical Kummer
construction; see \cite[Chapter~1, Example~1.3(iii)]{HuybrechtsK3Draft}.
Here the two required properties can also be checked in the
same charts. A translation-invariant two-form on $A$ is
invariant under $-1$. In a surface chart $a=r$, $b=ru$,
its pullback is
\[
                da\wedge db=r\,dr\wedge du
                            =\tfrac12\,dR\wedge du.   \tag{C.KM.7}\label{cdp:eq:C.KM.7}
\]
It therefore descends to a nowhere-zero form on $P$.
The finite covering $\operatorname{Bl}_{A[2]}A\to P$
injects $H^1(P,\OO_P)$ into the invariant part of
$H^1(A,\OO_A)$, by the averaging trace. This invariant
part is zero because $-1$ acts as $-1$ on $H^1(A,\OO_A)$.
Thus $K_P=\OO_P$ and $h^1(P,\OO_P)=0$, verifying that
$P$ is a K3 surface.

\noindent\textbf{Step 5.} The finite quotient map identifies rational cohomology with
invariants. Explicitly, its higher direct images of the
constant rational sheaf vanish by proper base change and
the finiteness of its fibres; averaging identifies the
constant sheaf downstairs with the invariant direct image.
Applied to $A/\{\pm1\}$, this gives Betti numbers
$(1,0,6,0,1)$ and Euler number eight. Resolving each of its
sixteen points replaces that point by $\PP^1$, so
$e(P)=8+16=24$. Hence the actual inclusion--exclusion
comparison in \eqref{cdp:eq:C.KM.1} gives
\[
          e\left(P\cup\bigcup_{i=1}^{16}E_i\right)
                    =24+16\cdot3-16\cdot2=40.          \tag{C.KM.8}\label{cdp:eq:C.KM.8}
\]
Mayer--Vietoris gives the support groups as well. The restriction
map from the direct sum of $H^2(P,\Q)$ and the sixteen
$H^2(E_i,\Q)$ onto the sixteen $H^2(P\cap E_i,\Q)$ is
surjective: the hyperplane class on each $E_i$ restricts to
the degree-one class of its intersection line. Its kernel
has dimension $22+16-16=22$, and the same exact sequence
gives $H^1(S_W,\Q)=H^3(S_W,\Q)=0$ and
$H^4(S_W,\Q)=\Q^{17}$. The coherent sequence
\[
 0\longrightarrow\OO_{S_W}\longrightarrow
 \OO_P\oplus\bigoplus_i\OO_{E_i}
 \longrightarrow\bigoplus_i\OO_{P\cap E_i}
 \longrightarrow0
\]
has a surjective difference map on $H^0$. The lines and planes
have no positive coherent cohomology, so it gives
$(h^0,h^1,h^2)(\OO_{S_W})=(1,0,1)$.

\noindent\textbf{Step 6.} We give the ambient cohomology maps as well. For the finite
quotient $A\times\PP^1\to X$, the same invariant-sheaf
argument gives
\[
 H^k(X,\Q)\simeq H^k(A\times\PP^1,\Q)^{\iota},
 \qquad (b_k(X))_{k=0}^6=(1,0,7,0,7,0,1).              \tag{C.KM.9}\label{cdp:eq:C.KM.9}
\]
Indeed $-1$ acts by $(-1)^j$ on $H^j(A)$, while the
base involution acts trivially on both nonzero cohomology
groups of $\PP^1$.

\noindent\textbf{Step 7.} Let $j_i:E_i\hookrightarrow W$ and let
$h_i=c_1(\OO_{E_i}(1))$. Collapsing the $32$ exceptional
divisors separately gives exactly $X$. The relative
cohomology sequence of this quotient identifies the
kernel of restriction to the $E_i$ with the pullback
of $H^k(X,\Q)$. The self-intersection formula gives
\[
 j_i^*j_{i*}(1)=-2h_i,
       \qquad j_i^*j_{i*}(h_i)=-2h_i^2.                 \tag{C.KM.10}\label{cdp:eq:C.KM.10}
\]
Thus restriction is surjective in degrees two and four,
with explicit rational sections $-j_{i*}(1)/2$ and
$-j_{i*}(h_i)/2$. In these two degrees pullback is
injective because the preceding odd cohomology of the
exceptional divisors vanishes. The same relative
sequence then leaves the other positive-degree groups
unchanged. Each exceptional point consequently adds one
class in degrees two and four and none in any other
degree. This proves the row for $W$ in \eqref{cdp:eq:C.KM.2}.

\noindent\textbf{Step 8.} For coherent cohomology the finite quotient and averaging
give
$H^k(X,\OO_X)=H^k(A\times\PP^1,\OO)^{\iota}$,
of dimensions $(1,0,1,0)$. The local resolution just
used has $R^q\pi_*\OO_W=0$ for $q>0$. To verify this
without suppressing the quotient singularity, filter the
infinitesimal neighbourhoods of $E_i$ by its ideal.
Their successive quotients are $\OO_{\PP^2}(2r)$,
$r\ge0$, whose positive-degree coherent cohomology
vanishes. Formal functions therefore makes the completed
stalks of those higher direct images zero, and coherence
makes the stalks themselves zero. Normality gives
$\pi_*\OO_W=\OO_X$. Leray proves the coherent row for
$W$.

\noindent\textbf{Step 9.} We now perform exactly twenty of the negative-two
operations of \ref{cdp:st:C.2.UN.4} over each of $0$ and $\infty$.
Every operation starts at a fresh smooth point on the
strict transform of one fixed $E_i$, away from its
intersection with $P$ and from all other components.
The first blowup creates a new $\PP^2$, and the second
blows up a smooth plane quartic of genus three inside
it, transverse to its old intersection line. These
centres and their entire modifications are disjoint
from $P$. The exact charts of \ref{cdp:st:C.2.UN.4} retain simple
normal crossings and show that both new components
have multiplicity one. Each operation changes the
reduced-fibre Euler number by $-2$ and adds two
components. Thus both final supports have Euler number
$40-20\cdot2=0$ and $17+20\cdot2=57$ components,
with their unchanged multiplicity-two K3 surfaces.

\noindent\textbf{Step 10.} We retain the support comparison under each of these blowups,
including the four boundary intersections of the quartic.
Let $\sigma:\widehat W\to W'$ be one blowup, with smooth
centre $C$ of ambient codimension $c$, let $S$ be the current
reduced support, and put $S'=\sigma^{-1}(S)_{\rm red}$ and
$E=\PP(N_{C/W'})$. Collapsing $E$ in $S'$ and $C$ in $S$
gives the actual homeomorphism $S'/E\simeq S/C$.
Hence pullback identifies the relative cohomology of the
pairs $(S,C)$ and $(S',E)$. With $p_E:E\to C$ and
$\xi=c_1(\OO_E(1))$, the projective-bundle decomposition is
\[
 H^k(E,\Q)=\bigoplus_{i=0}^{c-1}
                \xi^i p_E^*H^{k-2i}(C,\Q).
\]
For $i\ge1$, the restriction to $S'$ of the ambient Gysin
class $j_*(\xi^{i-1}p_E^*\alpha)$ restricts further to $E$
as $-\xi^i p_E^*\alpha$, because $N_{E/\widehat W}=\OO_E(-1)$.
These classes split every positive projective-bundle summand.
The two relative exact sequences now identify the remaining
summand with the pullback of $H^k(S,\Q)$ and give
\[
 H^k(S',\Q)\simeq H^k(S,\Q)\oplus
                \bigoplus_{i=1}^{c-1}H^{k-2i}(C,\Q).
                                                        \tag{C.KM.S2}\label{cdp:eq:C.KM.S2}
\]
This argument uses the full inverse image and does not require
the centre to avoid the other support components at every point.
A point step therefore adds one class in degrees two and four.
The quartic step adds ranks one, six and one in degrees two,
three and four. After twenty blocks, \eqref{cdp:eq:C.KM.S2} gives exactly
the second support row of \eqref{cdp:eq:C.KM.S1}.

\noindent\textbf{Step 11.} Each chosen centre is generically contained in just one
reduced component, of multiplicity one in the reduced support.
Consequently the equality of Cartier divisors
$\sigma^*S=S'$ holds, also at the four crossing points as
the explicit equations show. Apply derived pushforward to
$0\to\OO_{\widehat W}(-S')\to\OO_{\widehat W}\to
\OO_{S'}\to0$. The projection formula and
$R\sigma_*\OO_{\widehat W}=\OO_{W'}$ give
\[
                         R\sigma_*\OO_{S'}=\OO_S.       \tag{C.KM.S3}\label{cdp:eq:C.KM.S3}
\]
Thus all support coherent cohomology groups are unchanged,
proving the remaining assertion of \eqref{cdp:eq:C.KM.S1}.

\noindent\textbf{Step 12.} The pullback and exceptional Gysin maps in \ref{cdp:st:C.2.UN.4}
give cohomology increments $(\Delta b_2,\Delta b_3,
\Delta b_4)=(2,6,2)$ per operation. Forty operations
give the second Betti row in \eqref{cdp:eq:C.KM.2}. Smooth-centre
blowups preserve coherent cohomology, by their direct
image $\OO$ and vanishing higher direct images, so
the second coherent row is unchanged. This also checks
$e(Y)=0$ directly from its displayed Betti numbers.

\noindent\textbf{Step 13.} Finally, $P$ is not uniruled. A dominating ruled surface
would pull its nonzero holomorphic two-form back to a
nonzero two-form, whereas a ruled surface has no such
form. The inverse-map resolution in \ref{cdp:st:C.2.UN.3} therefore
prevents its contraction by any bimeromorphic comparison
between smooth threefold fillings. The finite-base-change
comparison is different and explicit: the map
\[
 A\times\PP^1_s\longrightarrow
 X\times_{\PP^1_t}\PP^1_s,
         \qquad(a,s)\longmapsto([a,s],s)               \tag{C.KM.11}\label{cdp:eq:C.KM.11}
\]
is finite and generically one-to-one. Its source is
normal, so it is the normalization of the reduced
base change. The normalized base changes of $W$ and
$Y$ factor properly and birationally through this
normalization. This specifies exactly how the Kummer
component is related to a torus after that base change.
The construction is based on the classical Kummer
quotient. Its chart,
trace, Euler and cohomology calculations are retained
in \path{code/E_kummer_filling.py} and its output.
\end{proof}
